% Part: incompleteness
% Chapter: incompleteness-provability
% Section: first-incompleteness-thm

\documentclass[../../../include/open-logic-section]{subfiles}

\begin{document}

\olfileid[id]{inc}{inp}{1in}

\olsection{Teorema Ketaklengkapan Pertama}

Sekarang kita dapat menguraikan pembuktian asli G\"odel atas teorema
ketaklengkapan pertama. Misalkan $\Th{T}$ sebarang teori yang diaksiomatisasi
secara komputabel dalam suatu bahasa yang memperluas bahasa aritmetika, dan
misalkan $\Th{T}$ memuat aksioma-aksioma~$\Th{Q}$. Ini berarti bahwa,
khususnya, $\Th{T}$ merepresentasikan fungsi dan relasi yang dapat dihitung.

Kita telah berargumen bahwa, dengan suatu pengodean formula dan pembuktian
yang wajar sebagai bilangan, relasi $\Prf[\Th{T}](x,y)$ dapat dihitung; di
sini $\Prf[\Th{T}](x,y)$ berlaku jika dan hanya jika $x$ adalah bilangan
G\"odel dari !!a{derivation} atas !!{formula} bernomor G\"odel~$y$
dalam~$\Th{T}$. Bahkan, untuk teori khusus yang dimaksud G\"odel, G\"odel
dapat menunjukkan bahwa relasi ini rekursif primitif dengan menggunakan
daftar 45 fungsi dan relasi dalam makalahnya. Relasi ke-45, $x B y$, tidak
lain adalah $\Prf[\Th{T}](x,y)$ untuk pilihan khusus~$\Th{T}$ miliknya.
Ingatlah bahwa ketika G\"odel menggunakan kata ``rekursif'' dalam
makalahnya, sekarang kita akan menggunakan frasa ``rekursif primitif.''

Karena $\Prf[\Th{T}](x,y)$ dapat dihitung, relasi itu dapat direpresentasikan
dalam~$\Th{T}$. Kita akan menggunakan $\OPrf[\Th{T}](x,y)$ untuk menyebut
formula yang merepresentasikannya. Misalkan $\OProv[\Th{T}](y)$ adalah
formula $\lexists[x][\OPrf[\Th{T}](x,y)]$. Formula ini menguraikan relasi
ke-46, $\fn{Bew}(y)$, dalam daftar G\"odel. Seperti dicatat G\"odel, inilah
satu-satunya relasi yang ``tidak dapat dinyatakan sebagai rekursif.'' Yang
mungkin ia maksud adalah ini: dari definisinya tidak jelas bahwa relasi itu
dapat dihitung; dan perkembangan selanjutnya memang menunjukkan bahwa relasi
itu tidak dapat dihitung.

Misalkan $\Th{T}$ suatu teori !!{axiomatizable} yang memuat~$\Th{Q}$. Maka
$\Prf[\Th{T}](x, y)$ dapat diputuskan, sehingga dapat direpresentasikan
dalam~$\Th{Q}$ oleh !!a{formula}~$\OPrf[\Th{T}](x, y)$. Misalkan
$\OProv[\Th{T}](y)$ adalah formula yang diuraikan di atas. Berdasarkan lema
titik tetap, terdapat suatu formula $!G_\Th{T}$ sedemikian sehingga
$\Th{Q}$ (dan karena itu $\Th{T}$) !!{derive}s
\begin{equation}
\ollabel{eqn:qpf}
!G_\Th{T} \liff \lnot \OProv[\Th{T}](\gn{!G_\Th{T}}).
\end{equation}
Perhatikan bahwa pada intinya $!G_\Th{T}$ mengatakan,
``$!G_\Th{T}$ tidak !!{derivable} dalam~$\Th{T}$.''

\begin{lem}\ollabel{lem:cons-G-unprov}
Jika $\Th{T}$ adalah teori konsisten dan !!{axiomatizable} yang
memperluas~$\Th{Q}$, maka $\Th{T} \Proves/ !G_\Th{T}$.
\end{lem}

\begin{proof}
Andaikan $\Th{T}$ !!{derive}s $!G_\Th{T}$. Maka \emph{memang} terdapat
!!a{derivation}, sehingga untuk suatu bilangan $m$, relasi
$\Prf[\Th{T}](m, \Gn{!G_\Th{T}})$ berlaku. Namun, kemudian $\Th{Q}$
!!{derive}s kalimat $\OPrf[\Th{T}](\num m, \gn{!G_\Th{T}})$. Jadi,
$\Th{Q}$ !!{derive}s $\lexists[x][\OPrf[\Th{T}](x,
\gn{!G_\Th{T}})]$, yang menurut definisi adalah
$\OProv[\Th{T}](\gn{!G_\Th{T}})$. Berdasarkan \olref{eqn:qpf}, $\Th{Q}$
!!{derive}s $\lnot !G_\Th{T}$; dan karena $\Th{T}$ memperluas $\Th{Q}$,
$\Th{T}$ juga !!{derive}s kalimat itu. Kita telah menunjukkan bahwa jika
$\Th{T}$ !!{derive}s $!G_\Th{T}$, maka teori itu juga !!{derive}s
$\lnot !G_\Th{T}$, sehingga teori itu akan tak konsisten.
\end{proof}

\begin{defn}
\ollabel{def:omega-consistency-inp}
Suatu teori $\Th{T}$ bersifat \emph{$\omega$-konsisten} jika hal berikut
berlaku: jika $\lexists[x][!A(x)]$ adalah sebarang kalimat dan $\Th{T}$
!!{derive}s $\lnot !A(\num 0)$, $\lnot !A(\num 1)$,
$\lnot !A(\num 2)$, \dots maka $\Th{T}$ tidak membuktikan
$\lexists[x][!A(x)]$.
\end{defn}

Perhatikan bahwa setiap teori $\omega$-konsisten juga konsisten. Hal ini
langsung mengikuti fakta bahwa jika $\Th{T}$ tak konsisten, maka
$\Th{T} \Proves !A$ untuk setiap~$!A$. Khususnya, jika $\Th{T}$ tak
konsisten, teori itu !!{derive}s baik $\lnot !A(\num n)$ untuk setiap~$n$
maupun~$\lexists[x][!A(x)]$. Jadi, jika $\Th{T}$ tak konsisten, teori itu
tidak $\omega$-konsisten. Berdasarkan kontraposisi, jika $\Th{T}$
$\omega$-konsisten, teori itu harus konsisten.

\begin{lem}\ollabel{lem:omega-cons-G-unref}
Jika $\Th{T}$ adalah teori $\omega$-konsisten dan !!{axiomatizable} yang
memperluas~$\Th{Q}$, maka $\Th{T} \Proves/ \lnot !G_\Th{T}$.
\end{lem}

\begin{proof}
Kita tunjukkan bahwa jika $\Th{T}$ !!{derive}s $\lnot !G_\Th{T}$, maka
teori itu tidak $\omega$-konsisten. Andaikan $\Th{T}$ !!{derive}s
$\lnot !G_\Th{T}$. Jika $\Th{T}$ tak konsisten, teori itu tak
$\omega$-konsisten, dan pembuktian selesai. Jika tidak, $\Th{T}$ konsisten,
sehingga teori itu tidak !!{derive} $!G_\Th{T}$ berdasarkan
\olref{lem:cons-G-unprov}. Karena tidak terdapat !!{derivation} atas $!G_\Th{T}$
dalam~$\Th{T}$, $\Th{Q}$ !!{derive}s
\[
\lnot \OPrf[\Th{T}](\num 0, \gn{!G_\Th{T}}), \lnot \OPrf[\Th{T}](\num 1,
\gn{!G_\Th{T}}), \lnot \OPrf[\Th{T}](\num 2, \gn{!G_\Th{T}}), \dots
\]
dan demikian pula~$\Th{T}$. Di sisi lain, berdasarkan \olref{eqn:qpf},
$\lnot !G_\Th{T}$ ekuivalen dengan
$\lexists[x][\OPrf[\Th{T}](x,\gn{!G_\Th{T}})]$. Jadi, $\Th{T}$ tak
$\omega$-konsisten.
\end{proof}

\begin{prob}
  Setiap teori $\omega$-konsisten bersifat konsisten. Tunjukkan bahwa
  konversnya tidak berlaku, yaitu bahwa terdapat teori yang konsisten tetapi
  tidak $\omega$-konsisten. Lakukan ini dengan menunjukkan bahwa $\Th{Q} \cup
  \{\lnot !G_\Th{Q}\}$ konsisten tetapi tidak $\omega$-konsisten.
\end{prob}

\begin{thm}
\ollabel{thm:first-incompleteness} Misalkan $\Th{T}$ sebarang teori
$\omega$-konsisten dan !!{axiomatizable} yang memperluas~$\Th{Q}$. Maka
$\Th{T}$ tidak lengkap.
\end{thm}

\begin{proof}
  Jika $\Th{T}$ $\omega$-konsisten, teori itu konsisten, sehingga
  $\Th{T} \Proves/ !G_\Th{T}$ berdasarkan
  \olref{lem:cons-G-unprov}. Berdasarkan
  \olref{lem:omega-cons-G-unref}, $\Th{T} \Proves/ \lnot !G_\Th{T}$.
  Ini berarti bahwa $\Th{T}$ tidak lengkap karena teori itu tidak !!{derive}s
  $!G_\Th{T}$ maupun $\lnot !G_\Th{T}$.
\end{proof}

\end{document}
